\chapter{Wykład I.}

\section{Wprowadzenie do równań różniczkowych.}

\subsection{Procent składany}
Rozważmy model procentu składanego. Oznaczenia
\begin{itemize}
	\item $u\left( 0 \right) $ -- stan rachunku w chwili początkowej,
	\item $r$-- oprocentowanie (np. w skali roku $r = 0.01$),
	\item $u\left( 1 \right) $ -- stan konta po jednym roku,
	\item $u\left( 1 \right) = u\left( 0 \right) + ru\left( 0 \right) $.
\end{itemize}
Mamy zatem mamy następujący model
\[
\begin{cases}
	u\left( 0 \right) &- \text{dane}\\
	u\left( n+1 \right) &= u\left( n \right) + ru\left( n \right)  
\end{cases}
.\] 
Rozwiązanie: $u\left( n \right) = u\left( 0 \right) \left( 1 + n \right)^{n}$.

\subsubsection*{Przejście do modelu ciągłego.}

Jeśli kapitalizacja następuje częściej (np. co pół roku, co chwilę), mamy
\begin{itemize}
	\item $u\left( 0 \right) $ -- dane,
	\item $u\left( \tfrac{1}{2} \right) = u\left( 0 \right) + \tfrac{1}{2}ru\left( 0 \right) $,
	\item $u\left( 1 \right) = u\left( \tfrac{1}{2} \right) + \tfrac{1}{2}u\left( \tfrac{1}{2} \right)$,
	\item itd
\end{itemize}

W modelu ciągłym wybieramy $h>0$ małe?, wtedy
\[
\begin{cases}
	h\left( 0 \right) &- \text{ dane,}  \\
	u\left( t+h \right) &= u\left( t \right) + hru\left( t \right)  
\end{cases}
.\] 
Równoważnie 
\[
\frac{u\left( t+h \right) - u\left( t \right) }{h} = ru\left( t \right) 
.\] 
W modelu ciągłym zakładamy, że $h \longrightarrow  0$, wtedy dostajemy \textbf{Prawo Malthusa}:
\[
\frac{du\left( t \right) }{dt} = ru\left( t \right) 
.\] 
To samo zjawisko jest w biologii
\begin{itemize}
	\item $r$ -- przyrost naturalny
	\item $u\left( t+h \right) = u\left( t \right) hrh\left( t \right) $
\end{itemize}

Matematycznie szukamy funkcji
\[
u = u\left( t \right) 
,\] 
która spełnia
\[
\begin{cases}
	\frac{du}{dt} &= rh \\
	u\left( 0 \right) &= u_{0} - \text{ dane } \\
\end{cases}
.\] 
Łatwo zgadnąć rozwiązanie: $u\left( t \right) = u_0e^{rt}$.
\subsection{Inne modele.}

\textbf{Model logistyczny:}

\[
\frac{du}{dt} = r\left( 1 - \frac{u}{L} \right)u
.\] 
$r>0$ oznacza przyrost, $r\in \mathbb{R}$. $L > 0$ oznacza optymalna pojemność populacji. Mamy
\begin{itemize}
	\item $r\left( 1 - \frac{u}{L} \right) > 0$, gdy $u < L$ 
	\item $r\left( 1 - \frac{u}{L} \right) > 0$, gdy $u > L$ 
\end{itemize}

\textbf{Model Gompertza}-- model rozprzestrzeniania się plotki
\[
frac{du}{dt} = ke ^{-lt}u
,\] 
gdzie $k>0$, $l>0$ -- stałe z doświadczenia.

\textbf{Prawo Newtona Stygnięcia}

$S\left( t \right) $ temperatura w danej chwili $t$.
\[
\frac{d}{dt}S = k\left( S_0 - S \right) 
,\]
gdzie $S_0$ jest to temperatura otoczenia.



